% Sefer Engine — Page ז
\definepapersize[seferpage][width=170mm,height=240mm]
\setuppapersize[seferpage]
\setuplayout[
  topspace=11mm, bottomspace=9mm, backspace=14mm,
  width=142mm, height=220mm, header=0mm, footer=0mm,
]
\mainlanguage[he]
\setupalign[r2l,hz,hanging]
\definefontfamily[hebrewfont][rm][David CLM]
\setupbodyfont[hebrewfont,12pt]
\setupindenting[no]
\setuppagenumbering[state=stop]

\starttext

% ═══ HEADER ═══
{\bf\tfx ז \hfill {\tfd שפע} \quad פרק א \hfill שלמה}
\blank[big]
{\bf ישראל ומשפיע עליהם כל טוב, ולא יחוס רק על כבוד הבורא ב״ה, כי זה כבודו יתברך שמו לראות בטובת ישראל", שכן על ידי טובתם מתגדל ומתקדש שמו יתברך בעולם [ב]. ולא עוד, אלא שכל השבח הגדול שמשבחים ומפארים מלאכי יתברך בעולם}
\blank[small]
\midaligned{\tfxx ◆ ◆ ◆ ◆ ◆ ◆ ◆ ◆ ◆ ◆ ◆ ◆ ◆ ◆ ◆ ◆ ◆ ◆ ◆ ◆ ◆ ◆ ◆ ◆}
\blank[small]
{\bf מקור השפע \hfill צינור השפע}
\blank[small]
% L-shape: makor longer → tzinor overlay LEFT, makor parshape RIGHT→FULL
\setbox0=\vbox{\hsize=69mm \setupalign[r2l,hz,hanging] \tfx\noindent [ב] וביסוד זה ביאר רבינו המבואר בגמרא (ברכות ז א) שכאשר נכנס רבי ישמעאל כהן גדול לפני ולפנים אמר לו הקב״ה "ישמעאל בני ברכני". וביאר רבינו: "כי הנה רבי ישמעאל הרהר איזה ברכה צריך הש״ת, הלא הוא מרומם על כל ברכה ותהלה? אך בודאי כוונת הש״ת הוא בשביל ישראל... שייטיב הש״ת עם ישראל עמו, ויהיה לכל אחד מישראל משאלות לבבו". \par\blank[small]\noindent באותה שעה שאמר רבינו דברים אלו הפנה את פניו לאנשים שעמדו סביבו והראה עליהם באצבעו: "זה צריך לפרנסה, וזה צריך לעשות שידוך עם בתו הבוגרת, וזה צריך לבן זכר, וזה צריך שיצליח במשפטו, וזה צריך למלבוש... וזאת היא הברכה אשר צריך לברך". והוסיף ואמר, כי אדם פשוט לא היה חושב ברוך יתברך ברכת הדיוט כמו פרנסה, אלא היה מברך יתגדל ויתקדש שמיה רבא וכדומה, אך רבי ישמעאל בן אלישע חשב במחשבתו שהש״ת אינו צריך להברכה לצורך עצמו, רק בעבור בני ישראל, שאם יהיה לישראל כל טוב זה הוא כבודו, כי הש״ת רוצה להיטיב לבני ישראל" (אהל שלמה ח״א אות א). \par\blank[small]\noindent בענין זה מסופר על הרה״ק משידלובצא זי״ע, שפעמו נחלם מאוד בין כסה לעשור, והרופאים הזהירו אותו שלא ילך למקום מפני הסכנה. בערב יום הכיפורים אמר שלא ישגיח על הרופאים וילך למקוה. אחד מתלמידיו טען לפניו שכמו שהכהן הגדול היה מתפלל תפילה קצרה "שלא להבעית את ישראל" (יומא נב ב), כך גם עליו להימנע ממעשה שיגרום לתלמידיו פחד ודאגה. על כך השיב הרבי בהתפעלות: "מה יקרו ומה נעמו דברי הרה״ק בעל תפארת שלמה מרדאמסק, שפירש כי זו עצמה היתה תפלת הכהן הגדול מהקב״ה: שלא להבעית את ישראל, כי ישראל יראים מיום הדין, יראים מן הגלות, יראים מן התורה הקדושה, יראים מעוה״ז ומעוה״ב, לכן אנא ה' עשה שלא להבעית את ישראל" (אהל שלמה ח״א אות כג).\par}
\vbox to 0pt{\hbox to \hsize{\hfill\copy0}\vss}%
\nointerlineskip
{\tfx
\parshape 114
  73mm 69mm
  73mm 69mm
  73mm 69mm
  73mm 69mm
  73mm 69mm
  73mm 69mm
  73mm 69mm
  73mm 69mm
  73mm 69mm
  73mm 69mm
  73mm 69mm
  73mm 69mm
  73mm 69mm
  73mm 69mm
  73mm 69mm
  73mm 69mm
  73mm 69mm
  73mm 69mm
  73mm 69mm
  73mm 69mm
  73mm 69mm
  73mm 69mm
  73mm 69mm
  73mm 69mm
  73mm 69mm
  73mm 69mm
  73mm 69mm
  73mm 69mm
  73mm 69mm
  73mm 69mm
  73mm 69mm
  73mm 69mm
  73mm 69mm
  73mm 69mm
  0mm 142mm
  0mm 142mm
  0mm 142mm
  0mm 142mm
  0mm 142mm
  0mm 142mm
  0mm 142mm
  0mm 142mm
  0mm 142mm
  0mm 142mm
  0mm 142mm
  0mm 142mm
  0mm 142mm
  0mm 142mm
  0mm 142mm
  0mm 142mm
  0mm 142mm
  0mm 142mm
  0mm 142mm
  0mm 142mm
  0mm 142mm
  0mm 142mm
  0mm 142mm
  0mm 142mm
  0mm 142mm
  0mm 142mm
  0mm 142mm
  0mm 142mm
  0mm 142mm
  0mm 142mm
  0mm 142mm
  0mm 142mm
  0mm 142mm
  0mm 142mm
  0mm 142mm
  0mm 142mm
  0mm 142mm
  0mm 142mm
  0mm 142mm
  0mm 142mm
  0mm 142mm
  0mm 142mm
  0mm 142mm
  0mm 142mm
  0mm 142mm
  0mm 142mm
  0mm 142mm
  0mm 142mm
  0mm 142mm
  0mm 142mm
  0mm 142mm
  0mm 142mm
  0mm 142mm
  0mm 142mm
  0mm 142mm
  0mm 142mm
  0mm 142mm
  0mm 142mm
  0mm 142mm
  0mm 142mm
  0mm 142mm
  0mm 142mm
  0mm 142mm
  0mm 142mm
  0mm 142mm
  0mm 142mm
  0mm 142mm
  0mm 142mm
  0mm 142mm
  0mm 142mm
  0mm 142mm
  0mm 142mm
  0mm 142mm
  0mm 142mm
  0mm 142mm
  0mm 142mm
  0mm 142mm
  0mm 142mm
  0mm 142mm
  0mm 142mm
\noindent
בפרשת תצוה (ד״ה א״י ואתה): "'בבני אהרן כתיב (ויקרא טז א) 'בקרבתם לפני ה' וגו', היינו שעבודתם היתה בלתי לה' לבדו, ולא היו בבחינה הנ״ל במסירות נפש לטובת ישראל, לכן זימותו'... כי שנים כאחד צריך להיות - כבוד הבורא ב״ה ולצורך טובת בני ישראל". \par\blank[small]\noindent בפרשת פנחס (ד״ה החתן): "יש שני מיני צדיקים, האחד, הוא המייחל ומצפה על גדולת הבורא ב״ה יהי כבוד ה' לעולם בביאת משיח במהרה בימינו, וזהו בחינת נדב ואביהוא - 'יקרבתם לפני ה'', שהיו נכנסים רק בבחינה זו לבדו ולא השגיחו כלל על טובת בני ישראל, לכן לא היה להם עוד חיות בעוה״ז, וזהו שנאמר 'בקרבתם לפני ה', [לפיכך] 'וימותו'. גם אצל אליהו כתיב (מלכים-א יט י) 'קנא קנאתי וגו' כי עוזבי בריתך וגו', וכתיב אח״כ (שם טו) 'ואת יואל אלישע בן שפט תמשח לנביא תחתיך". \par\blank[small]\noindent iט. פרשת שמיני (ד״ה אנא): "[הענין] במשה רבינו ע״ה שסירב בשליחות לגאול את בני ישראל ממצרים, הלא הוא היה אוהב נאמן ודורש טוב לעמו, ומדוע לא יחוש לגאלם. אך הנה ידוע כי גאולת מצרים לא היתה בשלימות מחמת כי דילג את הקץ, שהנגינה היתה ארבע מאות שנה, ואם השלימו העת ההוא היתה הגאולה בשלימות ולא היה שום שיעבוד וגלות עוד. וזהו שאמר לו ה' יתברך (שמות ג יד) 'אהיה אשר אהיה' באשר גליות. לכן סירב משה בשליחות הזה, כי רצונו היה יותר להתעכב עוד ולסבול צער הגלות עד שתמלא הקץ שיהיה הגאולה בשלימות שלא יהיה עוד גלות. ועיקר כוונת משרע״ה היה כדי שיתגדל ויתרומם כבודו יתברך שמו בעולם, כמו שאנו מצפים מהמהרה 'יהא שמיה רבא מברך' וכו'; ודאג. על כך השיב הרבי בהתפעלות: "מה יקרו ומה נעמו דברי הרה״ק בעל תפארת שלמה מרדאמסק, שפירש כי זו עצמה היתה תפלת הכהן הגדול מהקב״ה: שלא להבעית את ישראל, כי ישראל יראים מיום הדין, יראים מן הגלות, יראים מן התורה הקדושה, יראים מעוה״ז ומעוה״ב, לכן אנא ה' עשה שלא להבעית את ישראל" (אהל שלמה ח״א אות כג).
\par}

\stoptext
